% TOP_PROBLEM: {"id":30007538,"problem_number":"LOCAL-30007538","title":"Optimal sovereign-debt maturity","category_id":21,"category":{"id":21,"name":"economics","display_name":"Economics","description":"Open economic theory statements, published research agendas, and proposed scoped specifications; question provenance and historical priority are qualified per record.","slug":"economics","order_index":21,"created_at":"2026-10-08T00:00:00Z"},"status":"research_question","external_url":null,"legacy_ids":[],"published":false,"statement":"Find an optimal state-dependent sovereign maturity policy when term prices, default, and rollover risk are endogenous.","economics":{"original_id":27,"field":"Fiscal policy and sovereign debt","kind":"design","formalization_basis":"adapted_research_question","source_status":"explicit_open","priority_status":"earliest_found","priority_note":"This is the earliest explicit relevant statement verified in this audit, not a claim of original priority; older literature and earlier versions may contain antecedents.","earliest_verified_reference":{"authors":["Kris James Mitchener","Christoph Trebesch"],"title":"Sovereign Debt in the 21st Century","year":2021,"url":"https://www.ifo.de/DocDL/cesifo1_wp8959.pdf","doi":"10.3386/w28598","locator":"Section 7, p. 35, discussion of rollover risk and maturity management","evidence_summary":"The survey explicitly identifies measuring rollover crises and assessing policy responses, including debt maturity management and liquidity buffers, as unsolved challenges with scant empirical guidance."},"current_reference":{"authors":["Kris James Mitchener","Christoph Trebesch"],"title":"Sovereign Debt in the 21st Century","year":2021,"url":"https://www.ifo.de/DocDL/cesifo1_wp8959.pdf","doi":"10.3386/w28598","locator":"Section 7, p. 35, discussion of rollover risk and maturity management","evidence_summary":"The survey explicitly identifies measuring rollover crises and assessing policy responses, including debt maturity management and liquidity buffers, as unsolved challenges with scant empirical guidance."},"formalization_note":"The survey explicitly calls maturity-management assessment and rollover-risk measurement unresolved. The recursive optimization is a representative restricted formalization."},"statusQualification":"Published question or research agenda verified at the cited locator. The mathematical specification is adapted for this catalogue and includes additional assumptions; source publication does not establish first-ever priority or certify this exact model remains unsolved. The survey explicitly calls maturity-management assessment and rollover-risk measurement unresolved. The recursive optimization is a representative restricted formalization.","statusReviewedAt":"2026-10-08"}
% FIRST_POSTED: 2026-10-08
% ENTRY_KIND: statement_only
% !TeX program = lualatex
\documentclass[11pt]{article}
\usepackage{fontspec}
\usepackage[a4paper,margin=25mm]{geometry}
\usepackage{amsmath,amssymb,mathtools}
\usepackage{xurl}
\usepackage[hidelinks]{hyperref}
\newcommand{\sref}[2]{\hyperlink{source:#1:#2}{[S#2]}}
\setlength{\emergencystretch}{3em}
\begin{document}
% problemId:30007538
\section*{Optimal sovereign-debt maturity}\label{30007538}
\subsection{Definitions and mathematical statement}
Consider a sovereign with stochastic output \(y_t\), inherited short- and long-term debt \((b_t^S,b_t^L)\), and a specified cost of default. Short debt matures next period; long debt has declared coupon and decay rate. Investors price both claims using a stated discount factor, recovery rule, and information structure. Government decisions jointly determine primary surpluses, default, and new issuance; equilibrium prices \(q^S,q^L\) therefore depend on the policy and debt state.
Let \(m_t\in[0,1]\) be the fraction of new funding raised at long maturity. For a fixed initial state, determine a policy \(m^*(y,b^S,b^L)\) maximizing
\[
E\sum_{t\ge0}\beta^t
[u(c_t)-\mathcal D_t],
\]
where \(c_t\) is resident consumption after real default expenses, \(\beta\in(0,1)\) the resident discount factor, and \(\mathcal D_t\) additional nonconsumption disruption loss measured in utility units. Incorporate term premia, refinancing needs, and the opportunity to adjust fiscal policy after future information. Establish comparative statics and quantitatively disciplined welfare gains relative to fixed maturity shares. Separate fundamental insolvency from rollover failure through explicit equilibria or identification restrictions. Validate model-implied maturity and spread responses rather than fitting average debt composition alone. The target is optimal maturity in a declared endogenous-price environment; existing models and practical maturity strategies already solve important special cases, while empirical measurement and generalization remain unfinished.

\par\textbf{Formulation and scope.} The survey explicitly calls maturity-management assessment and rollover-risk measurement unresolved. The recursive optimization is a representative restricted formalization.

\par\textbf{Open-question provenance.} An explicit question or research agenda was verified at the source locator. This does not by itself certify that every later version remains unresolved \sref{30007538}{1}.

\par\textbf{Historical priority.} This is the earliest explicit relevant statement verified in this audit, not a claim of original priority; older literature and earlier versions may contain antecedents.

\subsection{Short English statement}
Find an optimal state-dependent sovereign maturity policy when term prices, default, and rollover risk are endogenous.

\subsection{Sources}
\begin{itemize}\raggedright
\item[\textnormal{[S1]}]\hypertarget{source:30007538:1}{} Kris James Mitchener, Christoph Trebesch. 2021. Sovereign Debt in the 21st Century. Section 7, p. 35, discussion of rollover risk and maturity management.
\url{https://www.ifo.de/DocDL/cesifo1_wp8959.pdf}.
DOI: 10.3386/w28598.
{\small\textit{Earliest verified question/agenda reference; historical priority qualified below. The survey explicitly identifies measuring rollover crises and assessing policy responses, including debt maturity management and liquidity buffers, as unsolved challenges with scant empirical guidance.}}
\end{itemize}
\end{document}
